% --- Radiacion Safety Manual Text in Spanish begins ---
\section*{Conceptos de Seguridad de Radiación}

También se usa el concepto ``capa de atenuación'' para indicar el espesor necesario para reducir la intensidad a un décimo de su valor original.

\subsection*{Cálculo de Capas Hemirreductoras}
El siguiente cuadro muestra los valores de las capas hemirreductoras para el concreto y el plomo y para algunos radioisótopos de interés:

% Table begins
\begin{table}[h]
\centering
\begin{tabular}{|c|c|c|c|}
\hline
\textbf{Fuente} & \textbf{Material} & \multicolumn{2}{c|}{\textbf{Espesor (pulgadas)}} \\
\cline{3-4}
 & & \textbf{Concreto} & \textbf{Plomo} \\
\hline
Iridio 192 & C.H. & 1.9" & .19" \\
 & C.D. & 6.2" & .64" \\
\hline
Cobalto 60 & C.H. & 2.6" & .49" \\
 & C.D. & 8.6" & 1.6" \\
\hline
Cesio 136 & C.H. & 2.1" & .25" \\
 & C.D. & 7.1" & .84" \\
\hline
\end{tabular}
\end{table}
% Table ends

Por ejemplo, si quisiéramos saber el espesor necesario de plomo para atenuar una fuente de Iridio-192 y reducir su tasa de exposición de 10 R/hr a 2 mR/hr, procederíamos aplicando la siguiente fórmula:
\[
X = \frac{X_0 }{ 2^n}
\]
donde \( X \) es la tasa de exposición, \( X_0 \) la tasa de exposición original y \( n \) el número de capas hemirreductoras que se necesitan.

% Equation for n
Como en este caso nos interesa conocer el \( n \), se despeja obteniendo la siguiente ecuación:
\[
n = \frac{\log(X_0 / X)}{\log 2}
\]
Sustituyendo los valores tenemos:
\[
n = \frac{\log(10,000/2)}{\log 2}
\]
\[
n = \frac{\log(5,000)}{\log 2}
\]
\[
n \approx 3.69897 / 0.30103 \approx 12.29
\]
Es decir, que se necesitan aproximadamente 12.29 capas hemirreductoras de plomo. Como cada una de ellas tiene 0.19", el espesor requerido será:
\[
\text{Espesor de plomo necesario} = 12.29 \times 0.19 = 2.34 \text{ pulgadas}.
\]
% --- Radiation Safety Manual Text in Spanish ends ---
